\documentclass[10pt,a4paper]{amsart}
\usepackage[margin=19mm]{geometry}
\usepackage{amsmath,amssymb,mathtools}
\usepackage[scr=rsfs,cal=euler]{mathalfa}
\usepackage{xfrac}
\usepackage[unicode,hidelinks,bookmarks=false]{hyperref}
\newcommand{\ie}{i.e.\ }
\newcommand{\cf}{cf.\ }
\newcommand{\resp}{respectively\ }
\newcommand{\Subsection}{\subsection}
\providecommand{\nobreakdash}{\nobreak\hbox{-}}
\setlength{\emergencystretch}{2em}
\multlinegap0pt
\usepackage{amssymb}
\usepackage{mathtools}
\usepackage[shortlabels]{enumitem}
\usepackage{xcolor}
\newtheorem{lemma}{Lemma}
\newtheorem{proposition}{Proposition}
\title{Entropy Expansion for General Polynomial Images\\ of Frostman Random Variables}
\author{Alex Iosevich, Thang Pham, Nguyen Dac Quan, Steven Senger, Boqing Xue}
\date{Source: arXiv:2507.15196v6 (CC BY 4.0)}
\begin{document}
\maketitle

\noindent\emph{Source and licence.} Entropy Expansion for General Polynomial Images\\ of Frostman Random Variables. Version arXiv:2507.15196v6 (CC BY 4.0), distributed by arXiv under the Creative Commons Attribution 4.0 International licence, http://creativecommons.org/licenses/by/4.0/, as recorded in the arXiv licence metadata for this exact version and shown on the abstract page https://arxiv.org/abs/2507.15196v6. Full licence notice: \texttt{licenses/arxiv-2507.15196-CC-BY-4.0.txt}.\par\smallskip

\noindent\textit{Editorial scope.} Counted unit: the article's proposition prop\_pham\_correlation (source0.tex lines 2375--2395). The article's complete original proof of it is delivered with it (source0.tex lines 2397--2501, one \texttt{proof} environment of 105 lines). No mathematics is written, repaired or re-derived: the statement and the proof are the article's own lines, in the article's own notation, and no retained line is substituted or re-flowed.  Retained uncounted as the source-local context the proof needs: lines 2106--2109; lines 2112--2115; lines 2131--2138; lines 2183--2207.  Nothing was omitted from any retained span, so the package carries no omission record.  No retained line contains a citation, so the wrapper carries no bibliography and no bibliography entry.  No retained line contains a figure environment, an image-inclusion command or drawing code, so no image is delivered and the package declares no asset dependency.  Editorial formatting only, in addition: an AMS \texttt{amsart} preamble that restates the article's own declarations and macros used by the retained lines, namely the environments \texttt{lemma}, \texttt{proposition} on separate counters, together with the standard packages those lines need.  The retained environments print in this document as Lemma 1, Proposition 1 in order of appearance, whereas the article prints the same items with the numbering of its own full document.  The complete native source of the article as distributed in the arXiv e-print for this exact version is bundled unchanged as \texttt{sources/181918/source0.tex}.
\begin{equation}
 I_a(\mu)\ll_{a,s,c}1+C.
 \label{eq_pham_input_energy}
\end{equation}

\begin{equation}
 H_n(\nu)\geq u n-\log I_u(\nu),\qquad 0<u<1.
 \label{eq_pham_energy_entropy}
\end{equation}

\begin{lemma}[Entropy and relative entropy]
\label{lem_pham_relative_entropy}
For any two probability vectors $p,q$,
\begin{equation}
 H(p)+2D(p\Vert q)\geq H_{\mathrm{coll}}(q).
 \label{eq_pham_relative_entropy}
\end{equation}
\end{lemma}

\begin{proposition}[Local folding energy estimate]
\label{prop_pham_local_energy}
Let $f\in\mathbb R[x,y]$ be non-special.  Let
\[
 K=I\times J\Subset K^+=I^+\times J^+
 \Subset\mathbb R^2\setminus Z(P_f),
\]
where $I,J,I^+,J^+$ are compact intervals.  Suppose that $\mu_1$ and
$\mu_2$ are probability measures supported on $I$ and $J$,
respectively.  If
\[
 0<a_1,a_2,u<1,
 \qquad a_1+a_2>u+\frac{2}{3},
\]
then
\begin{equation}
 I_u\big(f_*(\mu_1\times\mu_2)\big)
 \leq A
 \bigl(1+I_{a_1}(\mu_1)\bigr)
 \bigl(1+I_{a_2}(\mu_2)\bigr),
 \label{eq_pham_local_energy}
\end{equation}
where $A$ depends on the displayed fixed data, but not on the two
measures.
\end{proposition}

\begin{proposition}[Correlation penalty]
\label{prop_pham_correlation}
Let $f\in\mathbb R[x,y]$ be non-special.  Suppose that $(U,V)$ has
bounded range and is jointly $(a,b,K)$-Frostman, where
\[
 0<a,b\leq1,\qquad a+b>\frac{4}{3}.
\]
Put
\begin{equation}
 q(a,b)=\min\left\{1,a+b-\frac{2}{3}\right\}.
 \label{eq_pham_q_ab}
\end{equation}
For every $0<t<q(a,b)$, there are a fixed integer $r\geq0$ and a
constant $B<\infty$ such that
\begin{equation}
 H_n(f(U,V))+2\operatorname{I}_{n+r}(U;V)\geq t n-B
 \label{eq_pham_correlation}
\end{equation}
for every $n\geq1$.  The constants are uniform over all couplings
with the fixed displayed data.
\end{proposition}

\begin{proof}
Choose
\[
 t<u<q(a,b),\qquad a'<a,\quad b'<b,
 \qquad a'+b'>u+\frac{2}{3}.
\]
We first assume that the law $\mathsf P$ of $(U,V)$ is supported on one
product rectangle compactly contained in
$\mathbb R^2\setminus Z(P_f)$.  Let
$\mathsf Q=\mathsf P_U\times\mathsf P_V$.  Quantize the
two source coordinates at level $n+r$, where the fixed integer $r$
is chosen so that replacing each source cell by a representative
changes $f$ by $O(2^{-n})$.  Let $p$ and $q$ be the resulting
level-$n$ output laws under $\mathsf P$ and $\mathsf Q$, respectively.

The marginal laws are $a$- and $b$-Frostman.  Hence
\eqref{eq_pham_input_energy},
Proposition~\ref{prop_pham_local_energy}, and
\eqref{eq_pham_energy_entropy} give
\begin{equation}
 H_{\mathrm{coll}}(q)\geq u n-O(1).
 \label{eq_pham_product_collision}
\end{equation}
Replacing representatives by the original variables changes Shannon
entropy by $O(1)$.  It changes collision entropy by at most $O(1)$ as
well: bounded displacement ensures both that every old bin reaches
only $O(1)$ new bins and that every new bin receives mass from only
$O(1)$ old bins.  Cauchy--Schwarz therefore changes the squared
$\ell^2$ norm by at most a fixed factor.

Data processing for relative entropy gives
\begin{equation}
 D(p\Vert q)
 \leq D\big(\mathsf P_{n+r}\Vert
       (\mathsf P_U)_{n+r}\times(\mathsf P_V)_{n+r}\big)
 =\operatorname{I}_{n+r}(U;V).
 \label{eq_pham_data_processing}
\end{equation}
Lemma~\ref{lem_pham_relative_entropy},
\eqref{eq_pham_product_collision}, and the bounded displacement just
described imply
\begin{equation}
 H_n(f(U,V))+2\operatorname{I}_{n+r}(U;V)
 \geq u n-O(1)
 \label{eq_pham_local_correlation}
\end{equation}
on this rectangle.

It remains to localize uniformly.  A fixed real algebraic curve in a
bounded square has an $\eta$-neighborhood that can be covered by
$O(\eta^{-1})$ squares of side $O(\eta)$.  It follows from the joint
Frostman condition that
\begin{equation}
 \mathbb P\big(\operatorname{dist}((U,V),Z(P_f))<\eta\big)
 \ll \eta^{a+b-1}.
 \label{eq_pham_bad_tube}
\end{equation}
Fix $0<\delta<1-t/u$.  Choose $\eta>0$ so that the left-hand side of
\eqref{eq_pham_bad_tube} is at most $\delta$.  Partition the support
square into finitely many sufficiently small product cells.  Merge
the cells meeting the smaller bad neighborhood into one bad label;
each remaining cell has a fixed buffered enlargement disjoint from
$Z(P_f)$.  Let $L$ be the resulting finite label, let
$w_j=\mathbb P(L=j)$,
and call a label good when its cell is disjoint from the bad
neighborhood.  The total good mass $W$ satisfies
\begin{equation}
 W\geq1-\delta.
 \label{eq_pham_good_mass}
\end{equation}

Conditioning on a good cell of mass $w_j$ changes the two marginal
Frostman constants by at most a factor $O(1/w_j)$.  Applying the local
argument and \eqref{eq_pham_input_energy} therefore gives
\begin{align}
 &H(D_nf(U,V)\mid L=j)
 +2\operatorname{I}(D_{n+r}U;D_{n+r}V\mid L=j)\notag\\
 &\hspace{35mm}\geq u n-B_0-B_1\log\frac{1}{w_j},
 \label{eq_pham_cell_correlation}
\end{align}
where $B_0,B_1$ do not depend on $j$ or on the coupling.  The same
scale shift $r$ works for all good labels because there are only
finitely many cells.

Entropy concavity and conditional mutual information give
\begin{align*}
 \sum_{j\ \mathrm{good}}w_jH(D_nf(U,V)\mid L=j)
 &\leq H_n(f(U,V)),\\
 \sum_{j\ \mathrm{good}}w_j
 \operatorname{I}(D_{n+r}U;D_{n+r}V\mid L=j)
 &\leq\operatorname{I}(D_{n+r}U;D_{n+r}V\mid L)\\
 &\leq\operatorname{I}_{n+r}(U;V)+H(L).
\end{align*}
Moreover,
\[
 \sum_jw_j\log\frac{1}{w_j}=H(L)\leq\log|\operatorname{range}(L)|=O(1).
\]
Averaging \eqref{eq_pham_cell_correlation} over the good labels yields
\[
 H_n(f(U,V))+2\operatorname{I}_{n+r}(U;V)
 \geq Wu n-O(1)>t n-O(1),
\]
by the choice of $\delta$.  This proves
\eqref{eq_pham_correlation}.
\end{proof}

\end{document}
